\section{Related work}
\label{sec:related}

\paragraph{Autonomous research systems.} The systems literature we map is described in
Section~\ref{sec:results}; here we note only its shape. Surveys of the area exist
\citep{ramos2024llmchem}, but they are narrative and domain-scoped rather than systematic,
and predate the 2025--2026 expansion that accounts for most of our corpus. We are not aware
of a prior systematic map of this literature with a released protocol, coded corpus, and
reproducible pipeline.

\paragraph{Auditability and calibration of AI-produced claims.} A distinct and more recent
thread argues that the field's evaluation practices do not license its claims.
\citet{takahara2026hep} propose making the hypothesis--test--belief cycle an explicit,
auditable operation rather than an artifact of logs. \citet{wang2026firstresearch} apply the
same idea one stage earlier, requiring a research question to carry a certificate recording
its assumptions, mechanism, falsifier, and decisive test. \citet{li2026calibration} generalises
this into a framework of \emph{evidence-licensed claims}: validation does not determine claim
level, and automation amplifies the need for calibration. \citet{ning2026autoresearch}
supply the empirical counterpart, showing that agent-selected modelling changes can be
re-tested on data that never entered the search loop, and that most of them survive --- a
demonstration that held-out verification of agent decisions is practical. Our contribution to
this thread is measurement: we quantify, across a field-scale corpus, how often such
mechanisms are actually present.

\paragraph{Failure modes of automated research.} \citet{bowkis2026automated} argue that
automated alignment research will produce systematic errors concentrated where human
reviewers are least able to detect them, and that this is a property of optimisation rather
than of misbehaviour. \citet{ma2026sage} document the failure-recovery bottleneck in
autonomous experimenters and, notably, add a grounded-reporting mechanism that redacts
hallucinated numbers --- an admission that unconstrained agents fabricate results.
\citet{shi2026sdabench} and \citet{chen2026causalgame} independently find that current models
degrade sharply on inferential, mechanistic, and causal reasoning, with no evaluated agent
demonstrating reliable causal thinking.

\paragraph{Integrity of the scholarly record.} \citet{gyevnar2026ddos} show that indirect data
poisoning turns honest autonomous research agents into distributors of fraud, succeeding in
49.6\% of runs against three frontier systems while detected in 6\% --- and that a provenance
audit reduces the attack success rate to zero. \citet{brzozowski2026ghost} document the
downstream state of the record: 1{,}655 records on a DOI-minting repository authored by
non-existent people with model-family-specific name priors and server-verifiable backdating.
Both papers are in our corpus, and both bear on why the auditability numbers we report matter.

\paragraph{AI-conducted scholarship.} A small number of papers are themselves written with
declared AI participation, including autoethnographic work co-authored by the system under
study \citep{ying2026stratigraphy}. Our review differs in kind: rather than reporting on the
AI's introspection, we log the AI's externally checkable process --- decisions, failures,
and human interventions --- and analyse that log as data.

\section{Method}
\label{sec:method}

The protocol, codebook, screening prompts, and all amendments were written before the
corresponding stage and are released verbatim; amendments are dated and justified. We
summarise here and refer to the repository for the normative text.

\subsection{Search and corpus construction}

We harvested arXiv metadata via the public API using 42 query strings in three families:
system-identity terms (``AI scientist'', ``autonomous research agent'', ``self-driving
laboratory'', ``deep research'', \ldots), lifecycle-stage terms conditioned on agent or LLM
vocabulary, and audit/integrity terms. The search covered eight CS categories and the window
2022-01-01 to 2026-07-31, yielding \nHarvest{} unique records; \nScreened{} fell in the
codable window (first submission on or after 2024-01-01) or entered through the
expert-identified route. Per-query counts, including the absence of truncation, are recorded
in the released harvest manifest.

Two decisions in this stage are worth reporting because they were errors we caught. First, an
initial query conjoining ``auditable'' with agent vocabulary returned 2{,}336 records ---
68\% of the then-corpus --- because arXiv stems \texttt{auditable} to \texttt{audit*}, which
co-occurs with agent terminology throughout the safety literature. The query was replaced with
precise phrases and the harvest re-run. Second, a recall check against 15 landmark papers
known to the team found 13; the two misses used vocabulary outside every query family and
entered through the standard expert-identified route. Probing for recall also revealed that
our initial query set omitted ``deep research'' (343 in-window records), the term under which
much of the 2025--2026 work is published; supplementary queries were added before any
screening began.

A cross-source check against OpenAlex established a coverage limitation we report rather than
correct: flagship biological autonomous-science systems publish on bioRxiv, not arXiv. Our map
is of the arXiv-visible, computer-science-centric field.

\subsection{Screening}

Screening ran in two layers, both released.

\emph{Layer 1} screened every candidate twice, independently: Pass A on Claude Fable 5 with a
criteria-first prompt, Pass B on Claude Opus 5 with a question-first prompt. Using different
base models and different framings was intended to decorrelate rater error. Raw agreement was
\rawAgree{} ($\kappa = \kappaThree{}$ three-way; $\kappa = \kappaBinary{}$ for include versus
not), leaving \nAdjudicate{} decisions (\pctAdjudicate{}) requiring adjudication.

Inspecting those disagreements changed the review. They were not distributed randomly across
the corpus; they concentrated on seven recurring classes of paper --- listed in
Section~\ref{sec:reflexive} --- at the conceptual boundary of the field. The written criteria
simply did not decide them. We therefore wrote amendment A1: seven boundary rules with four
refined exclusion codes, resolving each class in the direction the research questions require
(the object is AI as a \emph{producer} of research, not AI as an aid to human researchers, not
detection of AI-written text, and not general-purpose agents).

Because the \nAgreedLayerOne{} agreed decisions had also been made without those rules,
applying them only to the disputed cases would have produced an internally inconsistent
corpus. \emph{Layer 2} therefore re-screened all \nScreened{} candidates under the amended
criteria in a single rule-application pass; it reversed \nReversals{} decisions on which both
Layer-1 passes had unanimously agreed. Layer 2 determines the corpus.

Amendment A2 records two decisions taken by the human co-author at a scope checkpoint. First,
the protocol's pre-specified narrowing at 400 papers was set aside as an artifact of a
workload guess rather than a scientific criterion; the corpus is reported in full. Second,
the human overruled boundary rule BR7, ruling general-purpose deep-research agents back into
scope on the grounds that excluding an entire product category because its benchmarks are
generic would define the field to fit the reviewer's framing. Re-screening the affected set
moved 183 papers into the corpus. We flag this as the clearest instance in the project of
human judgment correcting the AI's, and we report the AI's original position alongside it.

The final corpus is \nCorpus{} papers.

\subsection{Coding}

Each included paper was coded on ten dimensions: paper type; demonstrated autonomy level
(L0 assistive through L4 full autonomy); research-lifecycle stages performed; scientific
domain; evaluation methods; auditability mechanisms; code and data release; claim strength;
and the role assigned to humans. Coders were instructed to code what an abstract
\emph{demonstrates} rather than what it aspires to, and to mark a dimension \texttt{UNCLEAR}
rather than guess.

\subsection{Verifying artifact release from full text}
\label{sec:artifacts}

Coding artifact release from abstracts is biased, because repository links usually appear in
the body, a footnote, or a data-availability statement. Since this variable carries part of
our central claim, we measured it instead of inferring it: for every coded paper we downloaded
the PDF, extracted its text, and searched for repository and archive URLs with a fixed regular
expression. No model judgment is involved, so the procedure is deterministic and reproducible
from the released script.

The correction is large. Abstract-based coding put code release at \pctRepoAbstract{} of the
corpus; full-text verification finds repository links in \pctRepo{}. All artifact statistics
we report use the full-text measurement, and we flag the discrepancy as a caution for
abstract-based systematic reviews generally.

\subsection{Reliability, and what we did not do}

The protocol specified an independent human audit of a stratified sample of screening
decisions. The human co-author declined it, and we report the consequence plainly: no
screening decision in this review was independently validated by a human. In its place we ran
a third independent AI screening pass over a random sample of \nCons{} papers, which agreed
with the Layer-2 decision \consAgree{} of the time. That figure measures the stability of the
amended criteria under re-application, not their correctness: three passes sharing a model
family can be consistently wrong. We treat it as such throughout.
